% Cac muc cap nhat; thay tung muc tuong ung, khong dan ca tep vao cuoi bao cao.

\section{3.5. Thiết kế lịch sử thao
tác}\label{tv4-22-3-5-thiux1ebft-kux1ebf-lux1ecbch-sux1eed-thao-tuxe1c}

Mã TV3 của Dương Đình Phúc sử dụng ActionStack và UndoRedoManager. Mỗi
Action lưu type, startPos, endPos, content, previousContent,
cursorBefore và cursorAfter. Tọa độ dòng/cột của tầng lõi bắt đầu từ 0.
Nội dung cũ được giữ lại để khôi phục thao tác xóa hoặc thay thế.

ActionStack dùng danh sách liên kết kép: topNode giữ thao tác mới nhất,
bottomNode giữ thao tác cũ nhất. UndoRedoManager quản lý hai
ActionStack, không trực tiếp giữ hoặc chỉnh sửa DoublyLinkedList. Khi
Undo/Redo, bộ điều phối trả về Action để tầng gọi áp dụng lên tài liệu.

Khi mở tài liệu khác, onDocumentOpened xóa cả hai lịch sử. Trên giao
diện, cờ \_isApplyingUndoRedo ngăn việc ghi thao tác phục hồi thành
chỉnh sửa mới. Giới hạn k = 0 là không giới hạn; k \textgreater{} 0 giữ
tối đa k bản ghi Undo.


\section{4.11. Cài đặt lịch sử và kiểm tra ngoặc của
TV3}\label{tv4-22-4-11-cuxe0i-ux111ux1eb7t-lux1ecbch-sux1eed-vuxe0-kiux1ec3m-tra-ngoux1eb7c-cux1ee7a-tv3}

Phần này được đối chiếu với ActionStack, UndoRedoManager và
SyntaxChecker của Dương Đình Phúc trong bản mã
c7d3f205d9d67b14c676e14af8e119794da7236d.

\subsection{4.11.1. ActionStack và giới hạn k
bước}\label{tv4-22-4-11-actionstack-vuxe0-giux1edbi-hux1ea1n-k-bux1b0ux1edbc}

ActionNode lưu một Action cùng prev và next. push nối nút mới tại đỉnh;
pop lấy nút mới nhất; popBottom loại bỏ nút cũ nhất. Vì có con trỏ hai
đầu, cập nhật liên kết cho mỗi thao tác này là O(1), chưa tính sao chép
chuỗi trong Action. Khi push vượt k, nút ở đáy bị loại bỏ.
setMaxCapacity giảm giới hạn bằng cách loại các bản ghi cũ dư thừa; giá
trị âm bị từ chối, giá trị 0 chuyển sang không giới hạn. Hàm hủy, tạo
sao chép và gán sao chép quản lý các nút riêng để các Stack không cùng
sở hữu một chuỗi nút.

Ví dụ k = 3: thực hiện riêng lẻ a, b, c, d thì lịch sử còn b, c, d. Ba
lần Undo đưa tài liệu từ abcd về a; lần Undo thứ tư bị từ chối.

\subsection{4.11.2. Điều phối
Undo/Redo}\label{tv4-22-4-11-ux111iux1ec1u-phux1ed1i-undoredo}

recordAction kiểm tra tính hợp lệ, đưa Action vào undoStack và xóa
redoStack. Trích đoạn từ UndoRedoManager.cpp:

\begin{verbatim}
bool UndoRedoManager::recordAction(const Action& action) {
    if (!isValidAction(action)) return false;
    undoStack.push(action);
    redoStack.clear();
    return true;
}
\end{verbatim}

Undo lấy bản gốc từ undoStack, chuyển bản gốc sang redoStack và trả về
hành động nghịch đảo. Redo chuyển bản gốc trở lại undoStack rồi trả nó
cho tầng gọi thực thi. createInverseAction đổi chèn thành xóa và ngược
lại, đổi tách thành gộp dòng; thao tác thay thế hoán đổi content với
previousContent. cursorBefore và cursorAfter cũng được hoán đổi.

Ví dụ dán abc: Undo trả DELETE\_TEXT và tầng gọi xóa abc; Redo trả
INSERT\_TEXT và tầng gọi chèn lại abc. Nếu sau Undo người dùng dán xyz,
recordAction xóa nhánh Redo cũ. Bộ kiểm thử C12 áp dụng Action trả về
lên DoublyLinkedList rồi so sánh nội dung, không chỉ kiểm tra số bản
ghi.

\subsection{4.11.3. Kiểm tra ngoặc bằng
BracketStack}\label{tv4-22-4-11-kiux1ec3m-tra-ngoux1eb7c-bux1eb1ng-bracketstack}

SyntaxChecker dùng BracketStack liên kết đơn, tách biệt với ActionStack.
Mỗi BracketItem chứa ký tự ngoặc mở, dòng và cột. Ngoặc mở được push;
ngoặc đóng khớp đỉnh thì pop. Có ba loại lỗi: dư ngoặc đóng
(UNMATCHED\_CLOSING), sai cặp (MISMATCHED\_PAIR) và thiếu ngoặc đóng
(UNCLOSED\_OPENING). Ví dụ \texttt{({[}{]}\{\})} cân bằng,
\texttt{({[}){]}} sai thứ tự, \texttt{(()} thiếu \texttt{)}.

checkText kiểm tra chuỗi; checkDocument duyệt trực tiếp Node của
DoublyLinkedList. Khi ignoreStringsAndComments = true, bộ quét bỏ qua
chuỗi, hằng ký tự và chú thích //, /* \ldots{} */. BracketCheckResult
chứa trạng thái, loại lỗi, tọa độ, ký tự gặp/kỳ vọng và số ngoặc từng
loại. Thông báo hiển thị chuyển tọa độ từ 0-based sang 1-based. Nếu kết
thúc còn ngoặc mở và trước đó chưa ghi lỗi, hàm báo ngoặc còn trên đỉnh
Stack.

Với N byte và độ sâu ngoặc d, thời gian quét là O(N), bộ nhớ Stack O(d).
Cột được tính theo byte std::string, chưa phải vị trí ký tự Unicode. Bộ
quét không thay thế bộ phân tích cú pháp C++ đầy đủ; raw string và nối
dòng tiền xử lý chưa thuộc phạm vi kiểm thử hiện tại.


\section{5.3. Kiểm thử chức
năng}\label{tv4-22-5-3-kiux1ec3m-thux1eed-chux1ee9c-nux103ng}

\subsection{5.3.1. Kế hoạch kiểm thử toàn ứng
dụng}\label{tv4-22-5-3-kux1ebf-houx1ea1ch-kiux1ec3m-thux1eed-touxe0n-ux1ee9ng-dux1ee5ng}

TV4/KiemThu/KeHoachKiemThu.md chứa các ca thủ công TC01--TC17 cho bản
Windows đã tích hợp. Các ca C01--C14 dưới đây là kiểm thử C++ thuần,
không tự động xác nhận các nút, phím tắt hoặc giao diện Windows đã đạt.

\subsection{5.3.2. Kết quả kiểm thử tự động phần
lõi}\label{tv4-22-5-3-kux1ebft-quux1ea3-kiux1ec3m-thux1eed-tux1ef1-ux111ux1ed9ng-phux1ea7n-luxf5i}

KiemThuCore.cpp gọi DoublyLinkedList, ActionStack, UndoRedoManager và
SyntaxChecker. TV4 giữ C01--C10, mở rộng kiểm tra C11--C14 và xuất CSV.
Bản C01--C10 trước đó đã chạy trên Mac bằng clang++; bản mở rộng ngày
22/11/2026 được chạy lại trên Linux x86\_64 bằng GCC 14.2.0, C++17, -O2
-Wall -Wextra. Không gán lần chạy Linux này cho máy Mac.

\begingroup
\small
\setlength{\tabcolsep}{4pt}
\begin{longtable}{@{}p{0.07\linewidth}p{0.76\linewidth}p{0.12\linewidth}@{}}
\toprule
Mã & Nội dung kiểm thử & Kết quả \\
\midrule
\endhead
\bottomrule
\endlastfoot
C01 & Thêm, chèn, xóa dòng và liên kết hai chiều & PASS \\
C02 & Tách dòng rồi gộp lại & PASS \\
C03 & Chèn ký tự, Delete, Backspace & PASS \\
C04 & Chèn và xóa đoạn nhiều dòng & PASS \\
C05 & Sao chép, gán danh sách độc lập & PASS \\
C06 & Tìm phân biệt/không phân biệt hoa thường ASCII & PASS \\
C07 & Tìm các kết quả không chồng lấn & PASS \\
C08 & Thay tất cả và từ khóa không xuất hiện & PASS \\
C09 & Ranh giới dòng và xóa nút duy nhất & PASS \\
C10 & Lưu/mở lại văn bản UTF-8 & PASS \\
C11 & ActionStack: LIFO, lấy đáy, giảm k, sao chép, Stack rỗng & PASS \\
C12 & Undo/Redo trên tài liệu, xóa Redo, mở tệp mới, k = 3 & PASS \\
C13 & Ngoặc đúng/sai/thiếu/dư; chuỗi, chú thích và xuống dòng CRLF & PASS \\
C14 & Duyệt 50.000 dòng qua checkDocument, đếm ngoặc, thêm dòng lỗi & PASS \\
\end{longtable}
\endgroup

Tổng số: 14 ca; đạt: 14; không đạt: 0. Mỗi ca gồm một hoặc nhiều điều
kiện so sánh đầu ra. C14 kiểm tra dòng lỗi thứ 50.001 sau khi thêm một
dòng lỗi vào tài liệu 50.000 dòng.

CSV, log và môi trường chạy được lưu trong TV4/KetQua:
ket\_qua\_kiem\_thu\_core.csv, ket\_qua\_kiem\_thu\_core.log và
moi\_truong\_kiem\_thu\_core.txt. Cách chạy lại từ thư mục gốc dự án:

\begin{verbatim}
bash TV4/KiemThu/chay_kiem_thu_core.sh
\end{verbatim}

\subsection{5.3.3. Phạm vi kết
luận}\label{tv4-22-5-3-phux1ea1m-vi-kux1ebft-luux1eadn}

C01--C14 xác nhận các đầu ra đã kiểm tra của bốn lớp lõi. Bộ này chưa
chạy giao diện Windows, chưa phủ mọi trường hợp và chưa xác nhận các
phím tắt, ánh xạ vị trí Unicode hoặc lỗi I/O trên mọi hệ thống. C10 chỉ
kiểm tra giữ nội dung UTF-8 khi lưu/mở lại. Kết quả 14 ca là kiểm thử
chức năng; 120 lượt đo tại mục 5.7 vẫn là dữ liệu hiệu năng
DoublyLinkedList trên Mac và không phải phép đo Undo/Redo hay WinForms.


\section{5.4. Các bước chuẩn bị đã xác
nhận}\label{tv4-22-5-4-cuxe1c-bux1b0ux1edbc-chuux1ea9n-bux1ecb-ux111uxe3-xuxe1c-nhux1eadn}

\begin{longtable}[]{@{}
  >{\raggedright\arraybackslash}p{(\columnwidth - 4\tabcolsep) * \real{0.3333}}
  >{\raggedright\arraybackslash}p{(\columnwidth - 4\tabcolsep) * \real{0.3333}}
  >{\raggedright\arraybackslash}p{(\columnwidth - 4\tabcolsep) * \real{0.3333}}@{}}
\toprule\noalign{}
\begin{minipage}[b]{\linewidth}\raggedright
Nội dung
\end{minipage} & \begin{minipage}[b]{\linewidth}\raggedright
Bằng chứng
\end{minipage} & \begin{minipage}[b]{\linewidth}\raggedright
Phạm vi kết luận
\end{minipage} \\
\midrule\noalign{}
\endhead
\bottomrule\noalign{}
\endlastfoot
Sinh và đếm dữ liệu & Chương trình sinh dữ liệu, wc -l & Có ba tệp
10.000, 20.000, 50.000 dòng \\
Biên dịch bốn lớp C++ thuần & Script chay\_kiem\_thu\_core.sh & Các
mô-đun lõi biên dịch cùng nhau \\
Kiểm thử lõi mở rộng & CSV, log C01--C14 & 14/14 ca đạt trong môi trường
đã ghi \\
Chạy giao diện Windows & Chưa có bộ kết quả đầy đủ & Chưa kết luận
TC01--TC17 đã đạt \\
Đo thời gian quản lý văn bản & 120 lượt đo và 12 giá trị trung bình &
Giữ số liệu Mac ở mục 5.7--5.8 \\
\end{longtable}


\section{6.1. Kết quả hiện có}\label{tv4-22-6-1-kux1ebft-quux1ea3-hiux1ec7n-cuxf3}

Dự án đã có mã quản lý văn bản bằng DoublyLinkedList, lịch sử bằng
ActionStack và UndoRedoManager, cùng kiểm tra ngoặc bằng SyntaxChecker.
Phần giải thích tại mục 4.11 đã được đối chiếu trực tiếp với mã TV3.

TV4 đã chuẩn bị dữ liệu mẫu và ba bộ dữ liệu lớn, thu 120 lượt đo với 12
giá trị trung bình trên Mac. Bộ kiểm thử C++ thuần mở rộng đạt 14/14 ca
trong lần chạy Linux ngày 22/11/2026; kết quả và môi trường được lưu
cùng mã kiểm thử. Các bằng chứng này chưa thay thế việc kiểm thử toàn bộ
ứng dụng Windows đã tích hợp.


\section{6.2. Nội dung còn cần hoàn
thiện}\label{tv4-22-6-2-nux1ed9i-dung-cuxf2n-cux1ea7n-houxe0n-thiux1ec7n}

\begin{itemize}
\tightlist
\item
  Hoàn thiện kết nối nút Kiểm tra ngoặc với SyntaxChecker: trong bản
  c7d3f205, tsb\_check\_Click vẫn chưa có xử lý.
\item
  Chạy TC01--TC17 trên Windows, gồm Undo/Redo, giới hạn k và kiểm tra
  ngoặc.
\item
  Kiểm tra vị trí con trỏ và xử lý văn bản tiếng Việt qua giao diện.
\item
  Mở rộng kiểm thử lỗi đọc/ghi, trường hợp biên và mô hình cấp phát lỗi.
\end{itemize}

Không dùng nhãn thông báo đặt sẵn hoặc kết quả kiểm thử lõi để đánh dấu
toàn bộ ca giao diện là đạt.


\section{6.3. Hướng phát triển}\label{tv4-22-6-3-hux1b0ux1edbng-phuxe1t-triux1ec3n}

\begin{itemize}
\tightlist
\item
  Cải thiện ánh xạ chỉ số byte UTF-8 với vị trí ký tự trên giao diện.
\item
  Mở rộng bộ quét cho raw string và nối dòng tiền xử lý C++.
\item
  Tối ưu cập nhật và trực quan hóa khi tài liệu lớn.
\item
  Mở rộng kiểm thử lịch sử, con trỏ và dữ liệu mã nguồn thực tế.
\end{itemize}
